\chapter{球面幾何学}
    \section{余弦定理,正弦定理について}
        Wikipedia英語版の記事に簡潔な証明がある。
        3次元直角座標系の中心に単位球を置き、三角形の1頂点を北極に固定し、1つの辺を本初子午線に重ねてベクトルの内積を用いて余弦定理を導出している。
        その結果を用いて($\sin^2 A = 1 - \cos^2 A$)正弦定理が証明される。
    \section{極三角形の極三角形が元の三角形になること}
        \begin{proof}
            \quad\par
            $xyz$直角座標系を考え、原点を$O$とする。
            球の中心を$O$に据える。
            $A',B',C'$はそれぞれ$\overrightarrow{OB}\times\overrightarrow{OC},\;\overrightarrow{OC}\times\overrightarrow{OA},\;\overrightarrow{OA}\times\overrightarrow{OB}$と球面との交点である。
            よって$A,B,C$がそれぞれ$\overrightarrow{OB'}\times\overrightarrow{OC'},\;\overrightarrow{OC'}\times\overrightarrow{OA'},\;\overrightarrow{OA'}\times\overrightarrow{OB'}$の上にあることを示せば良い。
            すなわち、次の等式を示せば良い。
            \[
                \begin{cases}
                    \left[(\overrightarrow{OB}\times\overrightarrow{OC})\times(\overrightarrow{OC}\times\overrightarrow{OA})\right]\times\overrightarrow{OC} = \bm{0} \\
                    \left[(\overrightarrow{OC}\times\overrightarrow{OA})\times(\overrightarrow{OA}\times\overrightarrow{OB})\right]\times\overrightarrow{OA} = \bm{0} \\
                    \left[(\overrightarrow{OA}\times\overrightarrow{OB})\times(\overrightarrow{OB}\times\overrightarrow{OC})\right]\times\overrightarrow{OB} = \bm{0}
                \end{cases}
            \]
            1つ目の等式が成り立つことを示す。他の2つは同様にして示せる。
            ベクトル三重積の性質$\bm{a}\times(\bm{b}\times\bm{c}) = (\bm{a}\cdot\bm{c})\bm{b} - (\bm{a}\cdot\bm{b})\bm{c}$より
            \begin{align*}
                &[(\overrightarrow{OB}\times\overrightarrow{OC})\times(\overrightarrow{OC}\times\overrightarrow{OA})]\times\overrightarrow{OC} \\
                \quad &= \left\{\left[(\overrightarrow{OB}\times\overrightarrow{OC})\cdot\overrightarrow{OA}\right]\overrightarrow{OC} - \left[(\overrightarrow{OB}\times\overrightarrow{OC})\cdot\overrightarrow{OC}\right]\overrightarrow{OA}\right\}\times\overrightarrow{OC} = -\left\{\left[(\overrightarrow{OB}\times\overrightarrow{OC})\cdot\overrightarrow{OC}\right]\overrightarrow{OA}\right\}\times\overrightarrow{OC}\\
                \quad &= -\left\{\left[(\overrightarrow{OC}\times\overrightarrow{OC})\cdot\overrightarrow{OB}\right]\overrightarrow{OA}\right\}\times\overrightarrow{OC} \quad (\text{スカラー三重積の性質を用いた}) \\
                \quad &= \left[\left(\bm{0}\cdot{\overrightarrow{OB}}\right)\overrightarrow{OA}\right]\times\overrightarrow{OC} = \bm{0}
            \end{align*}
        \end{proof}
    \section{元の三角形\texorpdfstring{$ABC$}{ABC}の極三角形\texorpdfstring{$A'B'C'$}{A'B'C'}について\texorpdfstring{$A'=\pi-a,\;B'=\pi-b,\;C'=\pi-c$}{A'=π-a, B'=π-b, C'=π-c}となること(\texorpdfstring{$a=BC,\;b=CA,\;c=AB$}{a=BC, b=CA, c=AB})}
        \begin{proof}
            \quad\par
            $xyz$直角座標系を考え、原点を$O$とする。
            球の中心を$O$に据える。
            前述の定理より極三角形の極三角形は元の三角形だから、$C,B$はそれぞれ$\overrightarrow{OA'}\times\overrightarrow{OB'},\;\overrightarrow{OC'}\times\overrightarrow{OA'}$上にある。
            よって$\triangle OA'B'$と$\triangle OC'A'$の法線ベクトル同士の成す角が$\overrightarrow{OB}$と$\overrightarrow{OC}$の成す角すなわち$a$となる。
            $A'$は平面$OA'B'$と$OC'A'$の成す角(2つあり、その和は$\pi$)のうち$\triangle ABC$の辺$BC(=a)$に向かう側であるので、$A' = \pi-a$となる。
        \end{proof}